\documentclass[10pt,a4paper]{amsart}
\usepackage[margin=19mm]{geometry}
\usepackage{amsmath,amssymb,mathtools}
\usepackage[scr=rsfs,cal=euler]{mathalfa}
\usepackage{xfrac}
\usepackage[unicode,hidelinks,bookmarks=false]{hyperref}
\newcommand{\ie}{i.e.\ }
\newcommand{\cf}{cf.\ }
\newcommand{\resp}{respectively\ }
\newcommand{\Subsection}{\subsection}
\providecommand{\nobreakdash}{\nobreak\hbox{-}}
\setlength{\emergencystretch}{2em}
\multlinegap0pt
\usepackage{bm}
\usepackage{bbm}
\usepackage{dsfont}
\usepackage{xspace}
\usepackage{cancel}
\usepackage{mathtools}
\usepackage{amsthm}
\makeatletter
\@ifundefined{defn}{\newtheorem{defn}{Defn}}{}
\@ifundefined{thm}{\newtheorem{thm}[defn]{Theorem}}{}
\makeatother
\newcommand{\cB}{\mathcal{B}}
\title[The vanishing viscosity process for an eikonal equation in t]{The vanishing viscosity process for an eikonal equation in the radially symmetric setting}
\author{Fanchen Meng}
\date{Source: arXiv:2508.13230v1 (CC BY 4.0)}
\begin{document}
\maketitle

\noindent\emph{Source and licence.} The vanishing viscosity process for an eikonal equation in the radially symmetric setting. Version arXiv:2508.13230v1 (CC BY 4.0), distributed under the Creative Commons Attribution 4.0 International licence, as recorded in the arXiv licence metadata for this exact version and shown on the abstract page https://arxiv.org/abs/2508.13230v1. Full rights notice: licenses/arxiv-2508.13230-CC-BY-4.0.txt.\par\smallskip

\noindent\textit{Editorial scope.} Counted unit: the Thm whose source label is \texttt{None} in the article (source0.tex lines 551--553, source label \texttt{None}), delivered together with the article's own complete original proof of it (source0.tex lines 554--633, one \texttt{proof} environment of 80 lines). No mathematics is written, repaired or re-derived: statement and proof are the article's own text line for line. Required context retained uncounted: eikonal2 (lines 524--531). Every definition, display and label of those lines is retained unchanged. Omitted from this package and not redistributed: 1--523, 532--550, 634--953. Nothing inside a retained excerpt is omitted or altered: every retained body is delivered exactly as the article's lines are written, so no omission receipt of any kind accompanies this package. Citations: the retained excerpts carry no citation command at all, so no bibliography entry of the article is transcribed and no reference is redistributed. Reference adaptation: none was needed and none was made; every reference inside the delivered unit resolves inside it, so no unresolved reference or citation is produced. Printed slips kept literal: no printed word, symbol, hypothesis or number of the counted statement or of its proof was altered. Layout and class: the article's own class is \texttt{amsart} with packages including amsfonts, amsmath, amsthm, amssymb, epsfig, enumerate, bbm, fancyhdr, cleveref, geometry, graphicx, tikz, caption, color; the wrapper is the lane's pinned \texttt{amsart} editorial wrapper at 10pt on A4, which restates the theorem-like environments the retained text needs and adds the article's own macro definitions that the retained text needs (\texttt{\textbackslash cB}), with extra wrapper packages (bm, bbm, dsfont, xspace, cancel, mathtools). The delivered numbering is the unit's own. Assets: no retained span contains a figure environment, a graphics-inclusion command, a PDF-page inclusion, a table or a TikZ picture; no image file is copied into this package, the package declares no asset dependency and no image file is redistributed with it. Licence: version arXiv:2508.13230v1 of this article is distributed under the Creative Commons Attribution 4.0 International licence, as recorded in the arXiv licence metadata for this exact version and shown on the abstract page https://arxiv.org/abs/2508.13230v1. A full rights notice is delivered at licenses/arxiv-2508.13230-CC-BY-4.0.txt. Characters: every retained excerpt is pure ASCII, so the delivered engine, XeLaTeX with its default Latin Modern text font, reports no missing character.
\begin{equation}\label{eikonal2}
\left\{
\begin{aligned}
|Du| &= f(x) \quad \text{in } U,\\
u &= 0 \quad \text{on } \partial U,
\end{aligned} 
\right.    
\end{equation}\\

\begin{thm}
    Let $\mathcal{A}=\{x\in U| \;f(x)=0\}$ and $u$, $v$ be two viscosity solutions of \eqref{eikonal2}. Then, if $u\le v$ on $\mathcal{A}$, we have $u\le v$ in $U$.
\end{thm}

\begin{proof}
    Using the same notation as in the theorem. We suppose $u,v$ to be two viscosity solutions to \eqref{eikonal2} and $u\le v$ on $\mathcal{A}$.
    
    First, fix $\epsilon>0$. We consider the set $\mathcal{P}=\bigcup_{a\in\mathcal{A}}\cB(a,\delta)$, i.e., the union of balls with radius \( \delta \) centered at points in \( \mathcal{A} \). Since $u\le v$ in $\mathcal{A}$, we are able to choose $\delta_{\epsilon}>0$ small enough such that 
    \[
    su-(v+\epsilon)< -\frac{\epsilon}{2} \quad \text{in} \ \overline{\mathcal{P}}
    \]
for $s<1$ and $s$ close enough to 1.

For convenience, we denote $u^s = su$ and $v^{\epsilon}= v+\epsilon$ and our goal is to show $u^s \le v^{\epsilon}$ in $U$ for $\epsilon >0$ small and $s<1$ close to 1.\\
We continue our proof by using the classical ``doubling variable'' method. Assume by contradiction that 
\[
    \sup_{x\in \bar{U}\setminus\mathcal{P}}(u^s (x)-v^{\epsilon}(x)) = u^s (\bar x)-v^{\epsilon}(\bar x)= \sigma >0 \; \text{for some} \ \bar x\in \bar U\setminus \mathcal{P}.
\]
We see that:
\[
    u^{s}(x)-v^{\epsilon}(x) = su(x)-v(x)-\epsilon = 0-\epsilon<0 \; \text{on}\;\partial U \\
\]
and
\[
    u^s (x)-v^{\epsilon}(x) = su(x)-v(x)-\epsilon<-\frac{\epsilon}{2}<0 \;\text{on}\; \partial \mathcal{P}. 
\]\\
Hence, $\bar x $ can only belong to $U\setminus \mathcal{\bar P}$.

Here, we tried to avoid the case that $\bar x$ appears on the boundary because we will apply the viscosity subsolution and supersolution test later.
To be more precise, we need a local max or a local min near $\bar x$, which may not be true when $\bar x$ is on the boundary.
We consider an auxiliary function $\Psi(x,y) $:
\[
    \Psi^{\alpha}(x,y)=u^{s}(x)-v^{\epsilon}(y)-\frac{|x-y|^2}{2\alpha}
\]
and
\[
    \sup_{x,y\in \bar{U}\setminus\mathcal{P}}=\Psi ^{\alpha}(x_\alpha,y_\alpha) \; \text{for some} \; x_\alpha,y_\alpha \in \bar{U}\setminus\mathcal{P}.
\]
Since $(x_\alpha, y_\alpha) \rightarrow (\bar{x}, \bar{x})$ as $\alpha \rightarrow 0^+$, we have $x_\alpha, y_\alpha \in U \setminus \bar{\mathcal{P}}$ for sufficiently small $\alpha$.
Since $u$ is a viscosity solution to the \eqref{eikonal2}, we see that $u^s = su$ is a viscosity solution to the following equation:\\
\begin{equation}\label{eikonal3}
\left\{
\begin{aligned}
|Du^s| &= sf(x) \quad \text{in } U \\
u^s &= 0 \quad \text{on } \partial U
\end{aligned} 
\right.    
\end{equation}.\\
From the above, we know that the auxiliary function $\Psi^{\alpha}(x,y)$ has a max at $(x_\alpha,y_\alpha)$. Therefore, $x\mapsto \Psi^\alpha(x,y_\alpha)$ has a max at $x_\alpha$, which means
\[
    x\mapsto u^s (x)-\frac{|x-y_\alpha|^2}{2\alpha} \quad \text{has a max at} \; x_\alpha.
\]
As $u^s$ is a viscosity solution of \eqref{eikonal3}, by the viscosity subsolution test, we have 
\[
    \frac{1}{\alpha}|x_\alpha-y_\alpha|-sf(x_\alpha)\le0.
\]
Similarly, $y\mapsto \Psi^\alpha (x_\alpha,y)$ has a min at $y_\alpha$, which means 
\[
    y\mapsto v(y)-(-\frac{|x_\epsilon-y|^2}{2\alpha})\quad \text{has a min at} \; y_\alpha.
\]
As $v$ is a viscosity solution of \eqref{eikonal2}, by the viscosity supersolution test, we have 
\[
    \frac{1}{\alpha}|x_\alpha-y_\alpha|-f(y_\alpha)\ge0.
\]
Combine these two inequalities together, we get 
\[
    f(y_\alpha)\le \frac{1}{\alpha}|x_\alpha - y_\alpha|\le sf(x_\alpha).
\]
Take $\alpha \rightarrow 0^+$, we get 
\[
    f(\bar{x})\le sf(\bar{x}).
\]
Since $f(x)$ is nonnegative and  $\mathcal{A}=\{x\in U|f(x)=0\}$, then $f(x)>\gamma>0$ for some positive $\gamma$ in $U\setminus \bar{\mathcal{P}}$. Notice that  $\bar{x}\in U\setminus \bar{\mathcal{P}}$, $f(\bar{x})>0$. Recall that $s<1$, then the inequality 
\[
    f(\bar{x})\le sf(\bar{x})
\]
cannot be true. We get a contradiction. Hence we proved $u^s\le v^{\epsilon}$ in $U\setminus{\bar{\mathcal{P}}}$. 
Also, according to our construction at the beginning of the proof, we have $u^s\le v^\epsilon$ in $\mathcal{\bar{P}}$, too. As a result, $u^s\le v^\epsilon$ in the whole $U$. 

Moreover, $u^s\le v^\epsilon$ in $U$ is true for $s<1$ close to 1 and $\epsilon>0$ small. Then take $s\rightarrow 1^-$ and $\epsilon \rightarrow 0^+$, we have
\[
    u\le v \quad \text{in}\; U.
\]
\end{proof}

\end{document}
